% ========================================================
%  习近平新时代中国特色社会主义思想概论  期末模拟试卷（含参考答案）
%  覆盖范围：导论、第一章至第十二章（除第九章）
% ========================================================
\documentclass[12pt, a4paper]{article}

% ---- 中文支持 ----
\usepackage[UTF8]{ctex}

% ---- 数学 ----
\usepackage{amsmath, amssymb, amsthm}
\usepackage{mathtools}

% ---- 页面布局 ----
\usepackage[top=2.5cm, bottom=2.5cm, left=2.8cm, right=2.8cm]{geometry}

% ---- 表格 ----
\usepackage{booktabs, array, multirow, makecell}
\usepackage{tabularx}

% ---- 页眉页脚 ----
\usepackage{fancyhdr}
\usepackage{lastpage}
\pagestyle{fancy}
\fancyhf{}
\renewcommand{\headrulewidth}{0.6pt}
\lhead{\small 习近平新时代中国特色社会主义思想概论·期末模拟}
\rhead{\small 共 \pageref{LastPage} 页}
\cfoot{\small 第 \thepage\ 页}

% ---- 计数器与样式 ----
\usepackage{enumitem}
\usepackage{xcolor}
\usepackage{tcolorbox}
\tcbuselibrary{skins, breakable}

% ---- 填空线 ----
\newcommand{\blank}[1]{\underline{\hspace{#1}}}

% ---- 答案盒子 ----
\newtcolorbox{answerbox}[1][]{
  breakable,
  colback=gray!6,
  colframe=gray!50,
  fonttitle=\bfseries\small,
  title=参考答案,
  #1
}

% ---- 题号格式 ----
\newcounter{bigq}
\newcommand{\BigQ}[1]{%
  \stepcounter{bigq}%
  \vspace{6pt}%
  \noindent{\large\bfseries 第\chinese{bigq}大题\quad #1}%
  \vspace{4pt}\par
}

\newcounter{subq}[bigq]
\newcommand{\SubQ}{%
  \stepcounter{subq}%
  \noindent\textbf{（\arabic{subq}）}\quad
}

% 分隔线
\newcommand{\examrule}{\noindent\rule{\linewidth}{0.6pt}}

% ---- 文档开始 ----
\begin{document}

% ============================================================
%  封面
% ============================================================
\begin{center}
  {\LARGE\bfseries 习近平新时代中国特色社会主义思想概论}\\[8pt]
  {\large\bfseries 期末全真模拟试卷}\\[16pt]
  \begin{tabular}{lll}
    考试时间：120 分钟 & \quad & 满分：100 分 \\[4pt]
    \multicolumn{3}{l}{注意事项：请将答案写在答题纸指定区域，写在本卷上无效。}
  \end{tabular}
  \vspace{4pt}

  \examrule

  \vspace{6pt}
  \begin{tabular}{|c|c|c|c|c|c|}
    \hline
    \textbf{题号} & 一 & 二 & 三 & \textbf{总分} \\
    \hline
    \textbf{满分} & 30 & 20 & 50 & 100 \\
    \hline
    \textbf{得分} & & & & \\
    \hline
  \end{tabular}
\end{center}

\vspace{10pt}
\examrule
\vspace{6pt}

% ============================================================
%  第一大题：单项选择题
% ============================================================
\BigQ{单项选择题（每题 1 分，共 30 分）}

\SubQ 习近平新时代中国特色社会主义思想创立的时代背景不包括以下哪一项？

A. 世界百年未有之大变局加速演进

B. 中华民族伟大复兴进入关键时期

C. 中国式现代化全面推进拓展

D. 全球资本主义进入新发展阶段

\vspace{6pt}

\SubQ 以下哪项是习近平新时代中国特色社会主义思想最为核心关键的组成部分，起支撑作用的"四梁八柱"？

A. "十四个坚持"

B. "十三个方面成就"

C. "十个明确"

D. "六个必须坚持"

\vspace{6pt}

\SubQ "两个结合"中，马克思主义基本原理同中华优秀传统文化相结合，被比喻为：

A. "魂脉"与"根脉"的结合

B. "根脉"与"魂脉"的结合

C. "血脉"与"文脉"的结合

D. "主脉"与"支脉"的结合

\vspace{6pt}

\SubQ 习近平新时代中国特色社会主义思想的历史地位表述中，以下哪项是正确的？

A. 实现了马克思主义中国化时代化的"第三次飞跃"

B. 实现了马克思主义中国化时代化新的飞跃

C. 是马克思主义中国化的"最终成果"

D. 是"当代中国马克思主义、二十二世纪马克思主义"

\vspace{6pt}

\SubQ "两个确立"的决定性意义中，不包括以下哪项？

A. 确立习近平同志党中央的核心、全党的核心地位

B. 确立习近平新时代中国特色社会主义思想的指导地位

C. 确立中国特色社会主义制度的根本地位

D. 是新时代取得一切成就的根本原因

\vspace{6pt}

\SubQ 中国特色社会主义进入新时代，我国社会主要矛盾已经转化为：

A. 人民日益增长的物质文化需要同落后的社会生产之间的矛盾

B. 人民日益增长的美好生活需要和不平衡不充分的发展之间的矛盾

C. 人民对民主法治的需要同法治建设滞后之间的矛盾

D. 人民对生态环境的需要同环境污染严重之间的矛盾

\vspace{6pt}

\SubQ 关于我国社会主要矛盾的变化，以下表述正确的是：

A. 我国已经跨越社会主义初级阶段

B. 我国是世界最大发达国家的国际地位没有变

C. 我国仍处于并将长期处于社会主义初级阶段的基本国情没有变

D. 我国已经是中等发达国家

\vspace{6pt}

\SubQ 中国梦的本质是：

A. 国家强大、民族振兴、人民富裕

B. 国家富强、民族振兴、人民幸福

C. 国家统一、民族团结、人民安康

D. 国家现代化、民族复兴、人民小康

\vspace{6pt}

\SubQ 党的基本路线是国家的生命线、人民的幸福线，其核心内容可概括为：

A. "一个中心、两个基本点"

B. "两个中心、一个基本点"

C. "三个全面、四个坚持"

D. "五位一体、四个全面"

\vspace{6pt}

\SubQ 中国式现代化的中国特色中，以下哪项不是其特征？

A. 人口规模巨大的现代化

B. 全体人民同步富裕的现代化

C. 人与自然和谐共生的现代化

D. 走和平发展道路的现代化

\vspace{6pt}

\SubQ 推进中国式现代化必须牢牢把握的重大原则中，不包括：

A. 坚持和加强党的全面领导

B. 坚持以经济建设为中心

C. 坚持以人民为中心的发展思想

D. 坚持发扬斗争精神

\vspace{6pt}

\SubQ 中国最大的国情是：

A. 人口众多

B. 地域辽阔

C. 中国共产党的领导

D. 发展不平衡

\vspace{6pt}

\SubQ 中国共产党领导是中国特色社会主义制度的最大优势，主要是因为：

A. 党员人数最多

B. 党的自身优势是中国特色社会主义制度优势的主要来源

C. 党的历史最悠久

D. 党的组织最严密

\vspace{6pt}

\SubQ 维护党中央权威和集中统一领导，是：

A. 一般性的组织原则

B. 一个成熟的马克思主义执政党的重大建党原则

C. 临时性的政治要求

D. 仅适用于特殊时期的措施

\vspace{6pt}

\SubQ "江山就是人民，人民就是江山"的深刻内涵中，以下哪项表述最准确？

A. 人民是历史的创造者，是真正的英雄

B. 人民是国家的主人，享有最高权力

C. 人民是党的服务对象，需要被管理

D. 人民是经济发展的劳动力资源

\vspace{6pt}

\SubQ 人民立场是中国共产党的：

A. 基本政治立场

B. 根本政治立场

C. 重要政治立场

D. 一般政治立场

\vspace{6pt}

\SubQ 共同富裕是中国特色社会主义的：

A. 基本要求

B. 本质要求

C. 一般要求

D. 阶段性要求

\vspace{6pt}

\SubQ 习近平指出，"时代是出卷人，我们是答卷人，人民是阅卷人"，这体现了：

A. 人民是党的工作的最高裁决者和最终评判者

B. 党是最高政治领导力量

C. 时代决定一切

D. 考试制度的重要性

\vspace{6pt}

\SubQ 新时代全面深化改革开放就其艰巨性、复杂性和系统性来说，是：

A. 一场深刻的革命

B. 一次简单的政策调整

C. 一项常规的经济工作

D. 一个短期的改革任务

\vspace{6pt}

\SubQ 全面深化改革的总目标是：

A. 实现经济高速增长

B. 完善和发展中国特色社会主义制度、推进国家治理体系和治理能力现代化

C. 建立完全的市场经济体制

D. 实现全面私有化

\vspace{6pt}

\SubQ 新发展理念中，"协调发展"注重解决的是：

A. 发展动力问题

B. 发展不平衡问题

C. 人与自然和谐共生问题

D. 发展内外联动问题

\vspace{6pt}

\SubQ 构建新发展格局最本质的特征是：

A. 实现高水平的自立自强

B. 完全依赖国际市场

C. 放弃对外开放

D. 各地搞自我小循环

\vspace{6pt}

\SubQ 社会主义基本经济制度的新概括中，不包括以下哪项？

A. 公有制为主体、多种所有制经济共同发展

B. 按劳分配为主体、多种分配方式并存

C. 社会主义市场经济体制

D. 计划经济体制

\vspace{6pt}

\SubQ 教育、科技、人才是全面建设社会主义现代化国家的：

A. 一般性支撑

B. 基础性、战略性支撑

C. 临时性措施

D. 辅助性条件

\vspace{6pt}

\SubQ 人民民主是社会主义的：

A. 重要特征

B. 本质属性

C. 生命

D. 基本要求

\vspace{6pt}

\SubQ 全过程人民民主是社会主义民主政治的：

A. 一般形式

B. 本质属性

C. 补充形式

D. 临时措施

\vspace{6pt}

\SubQ 文化自信是"四个自信"中：

A. 最基础的自信

B. 更基础、更广泛、更深厚的自信

C. 最表面的自信

D. 最不重要的自信

\vspace{6pt}

\SubQ 意识形态关乎：

A. 经济发展

B. 旗帜、道路、国家政治安全

C. 文化建设

D. 国际关系

\vspace{6pt}

\SubQ 让人民生活幸福是：

A. "国之大者"

B. "党之大计"

C. "民之所盼"

D. "政之所向"

\vspace{6pt}

\SubQ 生态文明建设应该摆在：

A. 经济建设的次要位置

B. 全局工作的突出位置

C. 社会建设的补充位置

D. 文化建设的附属位置

\vspace{10pt}
\examrule

% ============================================================
%  第二大题：多项选择题
% ============================================================
\BigQ{多项选择题（每题 2 分，共 20 分）}

\SubQ 习近平新时代中国特色社会主义思想创立的时代背景包括：

A. 世界百年未有之大变局加速演进

B. 中华民族伟大复兴进入关键时期

C. 中国式现代化全面推进拓展

D. 科学社会主义在21世纪的中国焕发新的蓬勃生机

E. 中国共产党在革命性锻造中更加坚强有力

\vspace{6pt}

\SubQ "两个确立"的内容包括：

A. 确立习近平同志党中央的核心、全党的核心地位

B. 确立习近平新时代中国特色社会主义思想的指导地位

C. 确立中国特色社会主义制度的根本地位

D. 确立改革开放的基本国策

E. 确立以经济建设为中心的基本路线

\vspace{6pt}

\SubQ 新时代的内涵包括：

A. 承前启后、继往开来、在新的历史条件下继续夺取中国特色社会主义伟大胜利的时代

B. 决胜全面建成小康社会、进而全面建设社会主义现代化强国的时代

C. 全国各族人民团结奋斗、不断创造美好生活、逐步实现全体人民共同富裕的时代

D. 全体中华儿女勠力同心、奋力实现中华民族伟大复兴中国梦的时代

E. 我国不断为人类作出更大贡献的时代

\vspace{6pt}

\SubQ 推进中国式现代化需要正确处理的重大关系包括：

A. 顶层设计与实践探索的关系

B. 战略与策略的关系

C. 守正与创新的关系

D. 效率与公平的关系

E. 活力与秩序的关系

\vspace{6pt}

\SubQ 中国式现代化的本质要求包括：

A. 坚持中国共产党领导

B. 坚持中国特色社会主义

C. 实现高质量发展

D. 发展全过程人民民主

E. 实现全体人民共同富裕

\vspace{6pt}

\SubQ 党的自身优势包括：

A. 政治领导力

B. 思想引领力

C. 群众组织力

D. 社会号召力

E. 经济控制力

\vspace{6pt}

\SubQ 新发展理念的科学内涵包括：

A. 创新是引领发展的第一动力

B. 协调是持续健康发展的内在要求

C. 绿色是永续发展的必要条件

D. 开放是国家繁荣发展的必由之路

E. 共享是中国特色社会主义的本质要求

\vspace{6pt}

\SubQ 全过程人民民主的"四个相统一"包括：

A. 过程民主和成果民主的统一

B. 程序民主和实质民主的统一

C. 直接民主和间接民主的统一

D. 人民民主和国家意志的统一

E. 形式民主和内容民主的统一

\vspace{6pt}

\SubQ 中国特色社会主义政治制度的显著优势包括：

A. 坚持党的集中统一领导

B. 坚持人民当家作主

C. 坚持全面依法治国

D. 坚持全国一盘棋，集中力量办大事

E. 坚持各民族一律平等

\vspace{6pt}

\SubQ 建设社会主义文化强国需要：

A. 举旗帜、聚民心、育新人、兴文化、展形象

B. 坚持"二为"方向、"双百"方针

C. 激发全民族文化创新创造活力

D. 推动中华优秀传统文化创造性转化、创新性发展

E. 加强国际传播能力建设

\vspace{10pt}
\examrule

% ============================================================
%  第三大题：简答题
% ============================================================
\BigQ{简答题（每题 10 分，共 50 分）}

\SubQ（10 分）请阐述习近平新时代中国特色社会主义思想创立的时代背景，并说明"两个确立"的决定性意义。

\vspace{6pt}

\SubQ（10 分）为什么说中国特色社会主义进入新时代是我国发展新的历史方位？结合社会主要矛盾的变化，说明"一个变与两个没有变"的辩证关系。

\vspace{6pt}

\SubQ（10 分）请论述中国式现代化是如何创造人类文明新形态的，并说明推进中国式现代化必须牢牢把握的重大原则。

\vspace{6pt}

\SubQ（10 分）为什么说新时代全面深化改革开放是一场深刻革命？请结合全面深化改革的总目标和正确方法论加以论述。

\vspace{6pt}

\SubQ（10 分）请阐述新发展理念的科学内涵与实践要求，并说明如何以新发展理念引领高质量发展。

\vspace{10pt}
\examrule

\vspace{16pt}
\examrule
\vspace{4pt}
\begin{center}
  {\large\bfseries ——试题到此结束，请认真检查——}
\end{center}
\vspace{4pt}
\examrule

% ============================================================
%  参考答案
% ============================================================
\newpage
\setcounter{bigq}{0}

\begin{center}
  {\LARGE\bfseries 参考答案与评分标准}
\end{center}
\vspace{8pt}
\examrule

% -------- 第一大题参考答案 --------
\BigQ{单项选择题参考答案}

\begin{answerbox}

（1）D \quad （2）C \quad （3）A \quad （4）B \quad （5）C

（6）B \quad （7）C \quad （8）B \quad （9）A \quad （10）B

（11）B \quad （12）C \quad （13）B \quad （14）B \quad （15）A

（16）B \quad （17）B \quad （18）A \quad （19）A \quad （20）B

（21）B \quad （22）A \quad （23）D \quad （24）B \quad （25）C

（26）B \quad （27）B \quad （28）B \quad （29）A \quad （30）B

\end{answerbox}

% -------- 第二大题参考答案 --------
\BigQ{多项选择题参考答案}

\begin{answerbox}

（1）ABCDE \quad （2）AB \quad （3）ABCDE \quad （4）ABCDE \quad （5）ABCDE

（6）ABCD \quad （7）ABCDE \quad （8）ABCD \quad （9）ABCDE \quad （10）ABCDE

\end{answerbox}

% -------- 第三大题参考答案 --------
\BigQ{简答题参考答案}

\begin{answerbox}

\textbf{第1题（10分）：习近平新时代中国特色社会主义思想创立的时代背景与"两个确立"的决定性意义}

\textbf{一、时代背景（5分）}

习近平新时代中国特色社会主义思想创立的时代背景主要包括五个方面：

1. \textbf{世界百年未有之大变局加速演进}（1分）：国际力量对比深刻调整，全球治理体系加速变革，世界进入新的动荡变革期。

2. \textbf{中华民族伟大复兴进入关键时期}（1分）：我们比历史上任何时期都更接近、更有信心和能力实现中华民族伟大复兴的目标。

3. \textbf{中国式现代化全面推进拓展}（1分）：成功推进和拓展了中国式现代化，创造了人类文明新形态。

4. \textbf{科学社会主义在21世纪的中国焕发新的蓬勃生机}（1分）：中国的发展充分展示了科学社会主义的强大生命力。

5. \textbf{中国共产党在革命性锻造中更加坚强有力}（1分）：党经历深刻革命性锻造，自我净化、自我完善、自我革新、自我提高能力显著增强。

\textbf{二、"两个确立"的决定性意义（5分）}

"两个确立"是指：确立习近平同志党中央的核心、全党的核心地位；确立习近平新时代中国特色社会主义思想的指导地位。

其决定性意义在于：

1. "两个确立"是新时代取得一切成就的根本原因（1分），是党和国家事业取得历史性成就、发生历史性变革的最根本保证。

2. "两个确立"是战胜一切艰难险阻、应对一切不确定性的\textbf{最大确定性、最大底气、最大保证}（2分）。

3. "两个确立"反映了全党全军全国各族人民共同心愿，对新时代党和国家事业发展、对推进中华民族伟大复兴历史进程具有决定性意义（2分）。

\end{answerbox}

\vspace{6pt}

\begin{answerbox}

\textbf{第2题（10分）：新时代的历史方位与社会主要矛盾的辩证关系}

\textbf{一、中国特色社会主义进入新时代的历史方位（4分）}

中国特色社会主义进入新时代，是我国发展新的历史方位。这个新时代具有丰富内涵：

1. 是承前启后、继往开来、在新的历史条件下继续夺取中国特色社会主义伟大胜利的时代（1分）。

2. 是决胜全面建成小康社会、进而全面建设社会主义现代化强国的时代（1分）。

3. 是全国各族人民团结奋斗、不断创造美好生活、逐步实现全体人民共同富裕的时代（1分）。

4. 是全体中华儿女勠力同心、奋力实现中华民族伟大复兴中国梦的时代，也是我国不断为人类作出更大贡献的时代（1分）。

\textbf{二、"一个变与两个没有变"的辩证关系（6分）}

\textbf{1. "一个变"——社会主要矛盾的变化（2分）}

我国社会主要矛盾已经转化为\textbf{人民日益增长的美好生活需要和不平衡不充分的发展之间的矛盾}。这一变化反映了我国发展的阶段性特征，人民需求从"有没有"转向"好不好"，发展问题从"落后"转向"不平衡不充分"。

\textbf{2. "两个没有变"（2分）}

- 我国仍处于并将长期处于\textbf{社会主义初级阶段}的基本国情没有变。

- 我国是\textbf{世界最大发展中国家}的国际地位没有变。

\textbf{3. 辩证关系（2分）}

社会主要矛盾的变化是在社会主义初级阶段这个历史阶段中发生的变化，没有改变我们对我国社会主义所处历史阶段的判断。变与不变统一于中国特色社会主义伟大实践，既要看到变化带来的新机遇新挑战，又要立足基本国情、保持战略定力，坚持党的基本路线不动摇。

\end{answerbox}

\vspace{6pt}

\begin{answerbox}

\textbf{第3题（10分）：中国式现代化创造人类文明新形态与重大原则}

\textbf{一、中国式现代化创造人类文明新形态（5分）}

中国式现代化创造了人类文明新形态，主要体现在三个方面：

1. \textbf{提供全新现代化路径，超越西方现代化模式}（2分）：中国式现代化打破了"现代化=西方化"的迷思，证明现代化不是单选题，各国完全可以根据自身国情选择适合自己的现代化道路。它深深植根于中华优秀传统文化，体现科学社会主义的先进本质，借鉴吸收一切人类优秀文明成果，代表人类文明进步的发展方向。

2. \textbf{为广大发展中国家提供了新的选择}（2分）：中国式现代化展现了不同于西方现代化模式的新图景，为广大发展中国家独立自主迈向现代化树立了典范，提供了全新选择，为解决人类面临的共同问题提供了中国智慧、中国方案、中国力量。

3. \textbf{创造了人类文明新形态的理论和实践意义}（1分）：这一新形态是物质文明、政治文明、精神文明、社会文明、生态文明协调发展的文明形态，体现了社会主义建设规律和人类社会发展规律。

\textbf{二、推进中国式现代化必须牢牢把握的重大原则（5分）}

1. \textbf{坚持和加强党的全面领导}（1分）：这是中国式现代化的根本保证。

2. \textbf{坚持中国特色社会主义道路}（1分）：既不走封闭僵化的老路，也不走改旗易帜的邪路。

3. \textbf{坚持以人民为中心的发展思想}（1分）：让现代化建设成果更多更公平惠及全体人民。

4. \textbf{坚持深化改革开放}（1分）：不断增强社会主义现代化建设的动力和活力。

5. \textbf{坚持发扬斗争精神}（1分）：敢于斗争、善于斗争，知难而进、迎难而上。

\end{answerbox}

\vspace{6pt}

\begin{answerbox}

\textbf{第4题（10分）：新时代全面深化改革开放是一场深刻革命}

\textbf{一、为什么说新时代全面深化改革开放是一场深刻革命（4分）}

新时代全面深化改革开放，就其\textbf{艰巨性、复杂性和系统性}来说，是一场深刻的革命：

1. 这是\textbf{涉险滩、闯难关的变革}（1分），需要突破利益固化的藩篱，触及深层次体制机制问题。

2. 这是\textbf{思想理论的深刻变革、改革组织方式的深刻变革、国家制度和治理体系的深刻变革、人民广泛参与的深刻变革}（2分），涉及生产力与生产关系、经济基础与上层建筑的全面调整。

3. 这是在多年改革开放基础上的深化，是一场\textbf{全面、系统、整体的制度创新}（1分），使社会主义制度的优越性得到更好发挥。

\textbf{二、全面深化改革的总目标（3分）}

\textbf{完善和发展中国特色社会主义制度、推进国家治理体系和治理能力现代化}（1分）。

- "完善和发展中国特色社会主义制度"规定了改革的\textbf{根本方向}（1分）。

- "推进国家治理体系和治理能力现代化"明确了改革的\textbf{鲜明指向和时代要求}（1分）。

\textbf{三、全面深化改革开放要坚持的正确方法论（3分）}

1. \textbf{增强系统性、整体性、协同性}（1分）：坚持整体推进与重点突破相结合，促进各项改革举措在政策取向上相互配合、在实施过程中相互促进、在改革成效上相得益彰。

2. \textbf{加强顶层设计和摸着石头过河相结合}（1分）：鼓励勇于探索、勇于突破，不断形成新经验、产生新认识。

3. \textbf{统筹改革发展稳定}（1分）：坚持稳中求进工作总基调，把改革的力度、发展的速度和社会可承受的程度统一起来。

（此外，"胆子要大、步子要稳""坚持重大改革于法有据"等也是重要方法论要求，答出即可酌情给分。）

\end{answerbox}

\vspace{6pt}

\begin{answerbox}

\textbf{第5题（10分）：新发展理念的科学内涵与实践要求}

\textbf{一、新发展理念的科学内涵（5分）}

新发展理念包括\textbf{创新、协调、绿色、开放、共享}五大发展理念，分别回答不同根本问题：

1. \textbf{创新}是引领发展的第一动力，注重解决\textbf{发展动力问题}（1分）。

2. \textbf{协调}是持续健康发展的内在要求，注重解决\textbf{发展不平衡问题}（1分）。

3. \textbf{绿色}是永续发展的必要条件和人民对美好生活追求的重要体现，注重解决\textbf{人与自然和谐共生问题}（1分）。

4. \textbf{开放}是国家繁荣发展的必由之路，注重解决\textbf{发展内外联动问题}（1分）。

5. \textbf{共享}是中国特色社会主义的本质要求，注重解决\textbf{社会公平正义问题}（1分）。

\textbf{二、实践要求（3分）}

- 从\textbf{根本宗旨}上把握：坚持以人民为中心，把满足人民对美好生活的新期待作为发展的出发点和落脚点。

- 从\textbf{问题导向}上把握：切实解决发展不平衡不充分问题，解决影响构建新发展格局、影响人民群众生产生活的突出问题。

- 从\textbf{忧患意识}上把握：随时准备应对更加复杂困难的局面，牢牢守住安全发展底线。

\textbf{三、以新发展理念引领高质量发展（2分）}

高质量发展是能够很好满足人民日益增长的美好生活需要的发展，是体现新发展理念的发展。要更好统筹质的有效提升和量的合理增长，始终坚持\textbf{质量第一、效益优先}，推动经济发展从"有没有"转向"好不好"，为全面建设社会主义现代化国家提供更为坚实的物质基础。

\end{answerbox}

\vspace{12pt}
\examrule
\begin{center}
  \textit{参考答案仅供参考，评分时应酌情给分，观点正确、论述充分即可得分。}
\end{center}

\end{document}